%% 00_abstract.tex
\begin{abstract}
Clinical records are increasingly expressed in both structured standards (HL7 FHIR R4B) and
LLM-generated narrative summaries. When both representations encode the same facts, it is
unclear which should serve as the retrieval substrate for downstream question answering or
adherence-risk prediction.  We present a \emph{paired-data} empirical comparison that holds
information content constant: 200 synthetic patients rendered as FHIR~R4B bundles and as
LLM-generated narratives from the same underlying facts.  Three retrieval systems answer
13,800 programmatically-verified adherence-indicator questions across five Patient Support
Program (PSP) families (next-dose lookup, dose-history aggregation/MPR, coverage-window
reasoning/PDC, missed-dose detection, and persistence signalling) at three regimen-complexity
tiers.  A fourth arm extracts features from each representation for a synthetic adherence
classification task.

\textbf{Results.}  System~A (Narrative RAG) achieves 40.6\% exact-match accuracy
(95\% CI [39.8\%, 41.4\%]), outperforming System~B (Structured-Naive, 35.3\%) and
System~C (Structured-Aware, 33.4\%).  All pairwise gaps are statistically significant
(McNemar $p<0.001$; paired bootstrap CIs exclude zero).  An error taxonomy attributes
44--58\% of failures to reasoning truncation under the 512-token output budget; a
1024-token sensitivity sweep on a stratified 495-question sample recovers $+5.1$~pp
for System~C and $+3.8$~pp for System~B while leaving System~A unchanged, narrowing
but not eliminating A's lead.  Tier~3 multi-drug regimens collapse to $\approx$20\%
across all systems and remain so at the larger budget; this Tier-3 floor is below a
per-family most-frequent-class baseline at 24.8\%.  A no-retrieval baseline (System~N,
same answer LLM and prompt template, empty context) reaches 25.2\% overall and 6.9\%
at Tier~3, decomposing the A advantage into a $25.2\%$ question-text-attributable
floor plus a $+15.4$~pp retrieval-attributable lift; the corresponding lifts for B
and C are $+10.1$ and $+8.2$~pp.

\textbf{Feature extraction.}  On a synthetic adherence label whose construction is partly entailed by the
structured representation, FS-Structured features achieve AUC~0.997 against
FS-Narrative AUC~0.846 and FS-Aware AUC~0.769; we discuss the partial-tautology
in Results and Discussion.

\textbf{Contribution.}  We demonstrate and empirically characterise a \textit{representation
dissociation}: the representation that wins narrative QA (narrative $>$ structured) is not
the same representation that wins downstream ML adherence prediction (structured $\gg$
narrative).  A deterministic templated-narrative generator (System~A-T) eliminates LLM hallucination
risk and generation cost while matching or slightly exceeding the LLM narrative pipeline
on QA accuracy (41.3\% vs.\ 40.6\%; paired-bootstrap delta $+0.74$~pp [$+0.13$,
$+1.34$]); the narrative-format advantage therefore generalises beyond LLM-specific lexical
regularity.  To our knowledge, this is the first paired-data comparison of narrative RAG, naive
structured FHIR RAG, and resource-aware structured FHIR RAG in which information
content is held constant by shared-source rendering from a single bundle, and
ground truth is produced programmatically rather than by an LLM-as-judge.  Code (Apache-2.0)
and the synthetic dataset (CC-BY-4.0) are released as a combined Zenodo deposit
at \href{https://doi.org/10.5281/zenodo.20292792}{doi:10.5281/zenodo.20292792}.
\end{abstract}
